\documentclass[10pt]{article}
\usepackage[margin=0.75in]{geometry}
\usepackage{amsmath,microtype}
\setlength{\emergencystretch}{3em}
\title{Guide: Why a Fixed Odd-Step Budget Leaves a Gap}
\author{Collatz proof project}
\date{September 5, 2026}
\begin{document}
\maketitle
\section*{The proposed proof strategy}
The previous result reduces the escape problem to seeds whose multiplicative
coefficient never drops below one. One possible strategy is to assume a
smallest such seed exists, then find a smaller predecessor with the same
property.

Lean now proves a limitation of this strategy when the predecessor search
uses a fixed maximum number of odd steps. The limitation allows arbitrarily
many even steps; it is not just a short-search problem.

\section*{What the theorem says}
Fix an odd-step budget $K$. If a target has remainder one on division by
$3^K$, every positive-length segment reaching it with at most $K$ odd steps
has a contracting final coefficient.

Therefore no predecessor whose coefficients never contract can reach that
target within the budget. For example, a one-odd-step budget cannot supply
such a predecessor for any target congruent to one modulo three. More odd
steps can change the situation: $3\to5\to8\to4$ has two odd steps and
coefficient $9/8>1$.

\section*{Why it is true}
Reduce a starting seed modulo $2^t$, preserving the segment's parity word.
Lean proves that the endpoint of this canonical residue is below $3^j$,
where $j$ is its odd-step count. The original and canonical endpoints have
the same remainder modulo $3^j$.

If that remainder were one, the canonical endpoint would have to be one.
A noncontracting coefficient would then force the canonical seed to be one
as well. But a positive return from one to itself has a contracting
coefficient. This contradiction gives the obstruction.

\section*{What remains possible}
The obstructed targets can be arbitrarily large and can avoid multiples
of three. Thus adding a finite list of small exceptions cannot repair a
fixed-budget coverage claim.

None of those targets is proved to escape or to have coefficients that
never contract. The theorem identifies a gap in this particular backward
method. A budget that grows with the target, or another argument excluding
the residue class, remains possible. The noncontracting-seed exclusion and
nontrivial-cycle exclusion required for Collatz are still open.

\section*{Verification}
The Lean build and 144-declaration standard-foundations audit pass.
{\raggedright Run \texttt{lake build Collatz.InverseBarrier} from
\texttt{lean/}.\par}
The technical paper contains the exact statements and proof. No mathematical
priority claim is made.
\end{document}
